% !TeX root = ../main.tex
\section{задачи 3-3}

\vspace{1em}
\textbf{1. Теоретическое обоснование}

В системах с частотной модуляцией (ЧМ) увеличение индекса модуляции $\beta_{\text{чм}}$ позволяет повысить помехоустойчивость (выигрыш в отношении сигнал/шум на выходе детектора растет пропорционально $\beta_{\text{чм}}^2$. Однако это правило работает \textbf{только в надпороговом режиме}, когда мощность сигнала на входе детектора значительно превышает мощность шума в полосе пропускания высокочастотного тракта (БВЧ). 

Обычно пороговое отношение сигнал/шум по мощности на входе ЧМ-детектора принимается равным:
\begin{equation*}
    \rho_{\text{пор}} \approx 10
\end{equation*}
Если отношение сигнал/шум на входе падает ниже этого порога ($\rho_{\text{вх}} < 10$), наступает \textit{пороговый эффект}: полезный сигнал подавляется шумом, на выходе появляются аномальные выбросы (щелчки), и помехоустойчивость резко падает.

Отношение сигнал/шум по мощности на входе ЧД рассчитывается как:
\begin{equation*}
    \rho_{\text{вх}} = \frac{P_c}{P_{\text{ш}}} = \frac{P_c}{N_0 \cdot \Pi_{\text{БВЧ}}}
\end{equation*}
где $P_c = 2.5 \text{ мкВт} = 2.5 \cdot 10^{-6} \text{ Вт}$ — мощность сигнала, $N_0 = 10^{-12} \text{ Вт/Гц}$ — спектральная плотность мощности белого шума, а $\Pi_{\text{БВЧ}} = 2\Delta f$ — полоса пропускания БВЧ.

\vspace{1em}
\textbf{2. Расчет для исходного режима}

Исходный индекс модуляции: $\beta_{\text{чм1}} = 5$.
Частота модуляции: $F_{\text{м}} = 20 \text{ кГц} = 20 \cdot 10^3 \text{ Гц}$.

Девиация частоты:
\begin{equation*}
    \Delta f_1 = \beta_{\text{чм1}} \cdot F_{\text{м}} = 5 \cdot 20 \cdot 10^3 = 100 \text{ кГц} = 10^5 \text{ Гц}
\end{equation*}

Полоса пропускания БВЧ:
\begin{equation*}
    \Pi_{\text{БВЧ1}} = 2\Delta f_1 = 2 \cdot 10^5 \text{ Гц}
\end{equation*}

Мощность шума на входе ЧД:
\begin{equation*}
    P_{\text{ш1}} = N_0 \cdot \Pi_{\text{БВЧ1}} = 10^{-12} \cdot 2 \cdot 10^5 = 2 \cdot 10^{-7} \text{ Вт}
\end{equation*}

Отношение сигнал/шум на входе ЧД:
\begin{equation*}
    \rho_{\text{вх1}} = \frac{2.5 \cdot 10^{-6}}{2 \cdot 10^{-7}} = \frac{25}{2} = 12.5
\end{equation*}
Поскольку $\rho_{\text{вх1}} = 12.5 > 10$, в исходном состоянии приемник работает в линейном \textbf{надпороговом} режиме.

\vspace{1em}
\textbf{3. Расчет при увеличении индекса модуляции в 4 раза}

Новый индекс модуляции: $\beta_{\text{чм2}} = 4 \cdot 5 = 20$.

Новая девиация частоты:
\begin{equation*}
    \Delta f_2 = \beta_{\text{чм2}} \cdot F_{\text{м}} = 20 \cdot 20 \cdot 10^3 = 400 \text{ кГц} = 4 \cdot 10^5 \text{ Гц}
\end{equation*}

Новая требуемая полоса пропускания БВЧ:
\begin{equation*}
    \Pi_{\text{БВЧ2}} = 2\Delta f_2 = 2 \cdot 4 \cdot 10^5 = 8 \cdot 10^5 \text{ Гц}
\end{equation*}

Мощность шума в расширенной полосе:
\begin{equation*}
    P_{\text{ш2}} = N_0 \cdot \Pi_{\text{БВЧ2}} = 10^{-12} \cdot 8 \cdot 10^5 = 8 \cdot 10^{-7} \text{ Вт}
\end{equation*}

Новое отношение сигнал/шум на входе ЧД:
\begin{equation*}
    \rho_{\text{вх2}} = \frac{2.5 \cdot 10^{-6}}{8 \cdot 10^{-7}} = \frac{25}{8} = 3.125
\end{equation*}

\vspace{1em}
\textbf{4. Вывод}

При увеличении индекса модуляции в 4 раза отношение сигнал/шум на входе детектора упадет до значения $\rho_{\text{вх2}} = 3.125$. 
Так как это значение существенно меньше порога ЧМ-детектора ($\rho_{\text{вх2}} < 10$), наступит \textbf{пороговый эффект}. Полезный сигнал будет подавлен шумами, и качество приема катастрофически ухудшится.

\textbf{Ответ:} Нет, повысить помехоустойчивость таким образом нельзя. Расширение полосы пропускания приведет к падению отношения сигнал/шум на входе детектора ниже порогового значения, что вызовет резкое снижение помехоустойчивости приемника.